
\documentclass{article}
\usepackage{colortbl}
\usepackage{makecell}
\usepackage{multirow}
\usepackage{supertabular}

\begin{document}

\newcounter{utterance}

\centering \large Interaction Transcript for game `dond', experiment `coop\_en', episode 1 with Dialogue Architects\_\_qwen3.5{-}9b\_\_SCoRe\_Qwen3.5{-}9B.
\vspace{24pt}

{ \footnotesize  \setcounter{utterance}{1}
\setlength{\tabcolsep}{0pt}
\begin{supertabular}{c@{$\;$}|p{.15\linewidth}@{}p{.15\linewidth}p{.15\linewidth}p{.15\linewidth}p{.15\linewidth}p{.15\linewidth}}
    \# & $\;$A & \multicolumn{4}{c}{Game Master} & $\;\:$B\\
    \hline

    \theutterance \stepcounter{utterance}  
    & & \multicolumn{4}{p{0.6\linewidth}}{
        \cellcolor[rgb]{0.9,0.9,0.9}{
            \makecell[{{p{\linewidth}}}]{
                \texttt{\tiny{[P1$\langle$GM]}}
                \texttt{You are playing a cooperative negotiation game in which you have to agree on how to divide a set of items among you and another player.} \\
\\ 
\texttt{Rules:} \\
\texttt{(a) You and the other player are given a set of items. Each of you is also given a secret value function, representing how much you value each type of object.} \\
\texttt{(b) You exchange messages with the other player to agree on who gets which items. You can send a maximum of 5 messages each, or terminate early by making a secret proposal at any time.} \\
\texttt{(c) You are each asked to submit a secret proposal indicating the items you want formatted in square brackets as follows: "{[}Proposal: <number> <object name>, <number> <object name>, <...>{]}"} \\
\texttt{(d) If your proposals are complementary, i.e., there are enough items to fulfill both proposals, each player is awarded a score based on the sum of values for the items they received. Otherwise, both of you get zero points.} \\
\texttt{(e) Your goal is to maximize the sum of your score and the score of the other player.} \\
\\ 
\texttt{Let us start.} \\
\\ 
\texttt{The set of available items is:} \\
\\ 
\texttt{2 whistles, 2 mirrors, 2 keys} \\
\\ 
\texttt{Your secret value function is:} \\
\\ 
\texttt{whistle: 0, mirror: 2, key: 3} \\
\\ 
\texttt{IMPORTANT: Your messages, unless it is the secret proposal, are directly transmitted to the other player, so do not include any response to the rules or text announcing your message. To make a secret proposal, use the indicated format. Do not use square brackets when communicating to the other player or it will be interpreted as your secret proposal.} \\
            }
        }
    }
    & \\ \\

    \theutterance \stepcounter{utterance}  
    & \multicolumn{4}{p{0.6\linewidth}}{
        \cellcolor[rgb]{0.9,0.9,0.9}{
            \makecell[{{p{\linewidth}}}]{
                \texttt{\tiny{[P1$\rangle$GM]}}
                \texttt{I understand the rules and my valuation function. I am ready to negotiate.} \\
\\ 
\texttt{My valuations are:} \\
\texttt{{-} Whistle: 0} \\
\texttt{{-} Mirror: 2} \\
\texttt{{-} Key: 3} \\
\\ 
\texttt{Since I value whistles at 0, I should not ask for them. I value keys the most (3) and mirrors moderately (2). To maximize the total score, we need to ensure we don't both want the same high{-}value items, and we should aim for a split where the sum of our values is maximized.} \\
\\ 
\texttt{Let's start the negotiation.} \\
\\ 
\texttt{Hello! I see we have 2 whistles, 2 mirrors, and 2 keys. I value keys at 3, mirrors at 2, and whistles at 0. Since I don't care about whistles, I'd prefer to focus on keys and mirrors. How about we split the keys evenly (1 each) and split the mirrors evenly (1 each)? That way we both get a key and a mirror, and the whistles go to whoever wants them or are discarded. Does that sound fair to you?} \\
            }
        }
    }
    & & \\ \\

    \theutterance \stepcounter{utterance}  
    & & & \multicolumn{2}{p{0.3\linewidth}}{
        \cellcolor[rgb]{0.9,0.9,0.9}{
            \makecell[{{p{\linewidth}}}]{
                \texttt{\tiny{[GM$|$GM]}}
                \texttt{continue} \\
            }
        }
    }
    & & \\ \\

    \theutterance \stepcounter{utterance}  
    & & & \multicolumn{2}{p{0.3\linewidth}}{
        \cellcolor[rgb]{0.9,0.9,0.9}{
            \makecell[{{p{\linewidth}}}]{
                \texttt{\tiny{[GM$|$GM]}}
                \texttt{I understand the rules and my valuation function. I am ready to negotiate.} \\
\\ 
\texttt{My valuations are:} \\
\texttt{{-} Whistle: 0} \\
\texttt{{-} Mirror: 2} \\
\texttt{{-} Key: 3} \\
\\ 
\texttt{Since I value whistles at 0, I should not ask for them. I value keys the most (3) and mirrors moderately (2). To maximize the total score, we need to ensure we don't both want the same high{-}value items, and we should aim for a split where the sum of our values is maximized.} \\
\\ 
\texttt{Let's start the negotiation.} \\
\\ 
\texttt{Hello! I see we have 2 whistles, 2 mirrors, and 2 keys. I value keys at 3, mirrors at 2, and whistles at 0. Since I don't care about whistles, I'd prefer to focus on keys and mirrors. How about we split the keys evenly (1 each) and split the mirrors evenly (1 each)? That way we both get a key and a mirror, and the whistles go to whoever wants them or are discarded. Does that sound fair to you?} \\
            }
        }
    }
    & & \\ \\

    \theutterance \stepcounter{utterance}  
    & & \multicolumn{4}{p{0.6\linewidth}}{
        \cellcolor[rgb]{0.9,0.9,0.9}{
            \makecell[{{p{\linewidth}}}]{
                \texttt{\tiny{[GM$\rangle$P2]}}
                \texttt{You are playing a cooperative negotiation game in which you have to agree on how to divide a set of items among you and another player.} \\
\\ 
\texttt{Rules:} \\
\texttt{(a) You and the other player are given a set of items. Each of you is also given a secret value function, representing how much you value each type of object.} \\
\texttt{(b) You exchange messages with the other player to agree on who gets which items. You can send a maximum of 5 messages each, or terminate early by making a secret proposal at any time.} \\
\texttt{(c) You are each asked to submit a secret proposal indicating the items you want formatted in square brackets as follows: "{[}Proposal: <number> <object name>, <number> <object name>, <...>{]}"} \\
\texttt{(d) If your proposals are complementary, i.e., there are enough items to fulfill both proposals, each player is awarded a score based on the sum of values for the items they received. Otherwise, both of you get zero points.} \\
\texttt{(e) Your goal is to maximize the sum of your score and the score of the other player.} \\
\\ 
\texttt{Let us start.} \\
\\ 
\texttt{The set of available items is:} \\
\\ 
\texttt{2 whistles, 2 mirrors, 2 keys} \\
\\ 
\texttt{Your secret value function is:} \\
\\ 
\texttt{whistle: 2, mirror: 2, key: 1} \\
\\ 
\texttt{IMPORTANT: Your messages, unless it is the secret proposal, are directly transmitted to the other player, so do not include any response to the rules or text announcing your message. To make a secret proposal, use the indicated format. Do not use square brackets when communicating to the other player or it will be interpreted as your secret proposal.} \\
\\ 
\\ 
\texttt{I understand the rules and my valuation function. I am ready to negotiate.} \\
\\ 
\texttt{My valuations are:} \\
\texttt{{-} Whistle: 0} \\
\texttt{{-} Mirror: 2} \\
\texttt{{-} Key: 3} \\
\\ 
\texttt{Since I value whistles at 0, I should not ask for them. I value keys the most (3) and mirrors moderately (2). To maximize the total score, we need to ensure we don't both want the same high{-}value items, and we should aim for a split where the sum of our values is maximized.} \\
\\ 
\texttt{Let's start the negotiation.} \\
\\ 
\texttt{Hello! I see we have 2 whistles, 2 mirrors, and 2 keys. I value keys at 3, mirrors at 2, and whistles at 0. Since I don't care about whistles, I'd prefer to focus on keys and mirrors. How about we split the keys evenly (1 each) and split the mirrors evenly (1 each)? That way we both get a key and a mirror, and the whistles go to whoever wants them or are discarded. Does that sound fair to you?} \\
            }
        }
    }
    & \\ \\

    \theutterance \stepcounter{utterance}  
    & & & \multicolumn{4}{p{0.6\linewidth}}{
        \cellcolor[rgb]{0.9,0.9,0.9}{
            \makecell[{{p{\linewidth}}}]{
                \texttt{\tiny{[GM$\langle$P2]}}
                \texttt{That sounds like a solid starting point! Splitting the keys and mirrors evenly ensures we both get our highest{-}value items without conflict. Since I value whistles at 0, I'm happy to let them go to you if you want them, or we can just ignore them.} \\
\\ 
\texttt{To move forward, could you confirm your valuations for the items? Specifically, how much do you value a key, a mirror, and a whistle? Once I know that, I can calculate the total potential score and refine our split if needed.} \\
            }
        }
    }
     \\ \\

    \theutterance \stepcounter{utterance}  
    & & & \multicolumn{2}{p{0.3\linewidth}}{
        \cellcolor[rgb]{0.9,0.9,0.9}{
            \makecell[{{p{\linewidth}}}]{
                \texttt{\tiny{[GM$|$GM]}}
                \texttt{continue} \\
            }
        }
    }
    & & \\ \\

    \theutterance \stepcounter{utterance}  
    & & & \multicolumn{2}{p{0.3\linewidth}}{
        \cellcolor[rgb]{0.9,0.9,0.9}{
            \makecell[{{p{\linewidth}}}]{
                \texttt{\tiny{[GM$|$GM]}}
                \texttt{That sounds like a solid starting point! Splitting the keys and mirrors evenly ensures we both get our highest{-}value items without conflict. Since I value whistles at 0, I'm happy to let them go to you if you want them, or we can just ignore them.} \\
\\ 
\texttt{To move forward, could you confirm your valuations for the items? Specifically, how much do you value a key, a mirror, and a whistle? Once I know that, I can calculate the total potential score and refine our split if needed.} \\
            }
        }
    }
    & & \\ \\

    \theutterance \stepcounter{utterance}  
    & & \multicolumn{4}{p{0.6\linewidth}}{
        \cellcolor[rgb]{0.9,0.9,0.9}{
            \makecell[{{p{\linewidth}}}]{
                \texttt{\tiny{[P1$\langle$GM]}}
                \texttt{That sounds like a solid starting point! Splitting the keys and mirrors evenly ensures we both get our highest{-}value items without conflict. Since I value whistles at 0, I'm happy to let them go to you if you want them, or we can just ignore them.} \\
\\ 
\texttt{To move forward, could you confirm your valuations for the items? Specifically, how much do you value a key, a mirror, and a whistle? Once I know that, I can calculate the total potential score and refine our split if needed.} \\
            }
        }
    }
    & \\ \\

    \theutterance \stepcounter{utterance}  
    & \multicolumn{4}{p{0.6\linewidth}}{
        \cellcolor[rgb]{0.9,0.9,0.9}{
            \makecell[{{p{\linewidth}}}]{
                \texttt{\tiny{[P1$\rangle$GM]}}
                \texttt{Key: 3, Mirror: 2, Whistle: 0.} \\
\\ 
\texttt{Since we both value whistles at 0, they are essentially worthless to us. The total value of all items is:} \\
\texttt{{-} 2 Keys × 3 = 6} \\
\texttt{{-} 2 Mirrors × 2 = 4} \\
\texttt{{-} 2 Whistles × 0 = 0} \\
\texttt{Total Pool Value = 10.} \\
\\ 
\texttt{If we split the keys (1 each) and mirrors (1 each), we each get a value of 3 + 2 = 5. The sum of our scores would be 10, which is the maximum possible.} \\
\\ 
\texttt{I agree with your proposal: I take 1 key and 1 mirror. You take the other key and the other mirror. We can ignore the whistles.} \\
\\ 
\texttt{{[}Proposal: 1 key, 1 mirror{]}} \\
            }
        }
    }
    & & \\ \\

    \theutterance \stepcounter{utterance}  
    & & & \multicolumn{2}{p{0.3\linewidth}}{
        \cellcolor[rgb]{0.9,0.9,0.9}{
            \makecell[{{p{\linewidth}}}]{
                \texttt{\tiny{[GM$|$GM]}}
                \texttt{proposal} \\
            }
        }
    }
    & & \\ \\

    \theutterance \stepcounter{utterance}  
    & & & \multicolumn{2}{p{0.3\linewidth}}{
        \cellcolor[rgb]{0.9,0.9,0.9}{
            \makecell[{{p{\linewidth}}}]{
                \texttt{\tiny{[GM$|$GM]}}
                [0, 1, 1]
            }
        }
    }
    & & \\ \\

    \theutterance \stepcounter{utterance}  
    & & \multicolumn{4}{p{0.6\linewidth}}{
        \cellcolor[rgb]{0.9,0.9,0.9}{
            \makecell[{{p{\linewidth}}}]{
                \texttt{\tiny{[GM$\rangle$P2]}}
                \texttt{The other player has submitted a secret proposal. It is now your turn to submit a proposal of your own.} \\
            }
        }
    }
    & \\ \\

    \theutterance \stepcounter{utterance}  
    & & & \multicolumn{4}{p{0.6\linewidth}}{
        \cellcolor[rgb]{0.9,0.9,0.9}{
            \makecell[{{p{\linewidth}}}]{
                \texttt{\tiny{[GM$\langle$P2]}}
                \texttt{{[}Proposal: 1 key, 1 mirror{]}} \\
            }
        }
    }
     \\ \\

    \theutterance \stepcounter{utterance}  
    & & & \multicolumn{2}{p{0.3\linewidth}}{
        \cellcolor[rgb]{0.9,0.9,0.9}{
            \makecell[{{p{\linewidth}}}]{
                \texttt{\tiny{[GM$|$GM]}}
                \texttt{proposal} \\
            }
        }
    }
    & & \\ \\

    \theutterance \stepcounter{utterance}  
    & & & \multicolumn{2}{p{0.3\linewidth}}{
        \cellcolor[rgb]{0.9,0.9,0.9}{
            \makecell[{{p{\linewidth}}}]{
                \texttt{\tiny{[GM$|$GM]}}
                [0, 1, 1]
            }
        }
    }
    & & \\ \\

    \theutterance \stepcounter{utterance}  
    & & & \multicolumn{2}{p{0.3\linewidth}}{
        \cellcolor[rgb]{0.9,0.9,0.9}{
            \makecell[{{p{\linewidth}}}]{
                \texttt{\tiny{[GM$|$GM]}}
                [[0, 1, 1], [0, 1, 1]]
            }
        }
    }
    & & \\ \\

\end{supertabular}
}

\end{document}
